% Fraction d'acide hypochloreux en fonction du pH (exercice 15, chapitre 9)
\begin{tikzpicture}
\begin{axis}[otMesures,otAxesBout,otLecture,width=100mm,height=50mm,xmin=5.4,xmax=9.6,ymin=0,ymax=1.08,
  xtick={5.5,6,...,9.5},ytick={0,0.1,...,1},
  xticklabel style={/pgf/number format/fixed,/pgf/number format/precision=1},
  yticklabel style={/pgf/number format/fixed,/pgf/number format/precision=1},
  xlabel={pH},ylabel={fraction de \ce{HClO}}]
\fill[otGreenTint] (axis cs:7.2,0) rectangle (axis cs:7.4,1.0);
\node[pcLegende,text=otGreenDark,anchor=south,align=center] at (axis cs:7.3,1.0) {plage\\recommandée};
\addplot[otNavy,line width=1pt,domain=5.4:9.5,samples=160] {1/(1+10^(x-7.5))};
\solution{\draw[otOrange,line width=0.6pt,dash pattern=on 2pt off 1.5pt]
  (axis cs:7.2,0) -- (axis cs:7.2,0.666) -- (axis cs:5.4,0.666)
  (axis cs:8.2,0) -- (axis cs:8.2,0.166) -- (axis cs:5.4,0.166);}
\end{axis}
\end{tikzpicture}
\par\smallskip\legende{Fraction du chlore sous forme d'acide hypochloreux, la forme désinfectante, en fonction du pH de l'eau.}
